%% ==========================================================================
%%  6  What is Proved
%%
%%  MAIN REPORT ONLY -- not included by working.tex, which carries the full
%%  development.  Every label here is prefixed `main:' so nothing can collide
%%  with working/.
%%
%%  Condensed from working/solved_cases.tex (thm:binadd, thm:identical,
%%  thm:smallbundle, cor:m-le-n, thm:balanced-class) and
%%  working/approach_10.tex (thm:f5-sigma3 and the reduction around it).
%% ==========================================================================
\section{What is Proved}
\label{sec:results}

Conjecture~\ref{conj:main} is open. Several restrictions of it are not, and they
are collected here because together they delimit what a general proof still has
to reach. The first four are cases of the valuation class or of the instance.
The fifth is a criterion at three agents, and it yields the last: the first case
of the conjecture proved for any $n \ge 3$.

% --------------------------------------------------------------------------
\subsection{Four solved cases}
\label{sec:main-solved}

\begin{theorem}[Binary additive costs]
\label{main:binadd}
Conjecture~\ref{conj:main} holds when every $\cost_i$ is additive with all
singleton values in $\set{0,1}$, that is
$\cost_i(S) = \abs{S \cap D_i}$ for some $D_i \subseteq \items$.
\end{theorem}

This is~\cite{LMS26}; the contribution here is only the observation that the
class sits inside the dichotomous one.

\begin{theorem}[Identical costs]
\label{main:identical}
If $\cost_1 = \dots = \cost_n = \cost$ then every allocation is envy-freeable,
and greedily adding each chore to a currently cheapest bundle yields
$\optsubsidy \in \set{0,1}^n$.
\end{theorem}

\begin{proof}
For any permutation $\sigma$, $\sum_i \cost(A_{\sigma(i)}) = \sum_i
\cost(\bundle{i})$, so condition~(ii) of Theorem~\ref{thm:hs-characterisation}
holds with equality and every allocation is envy-freeable. Moreover
$\edgew{\alloc}(i,j) = \cost(\bundle{i}) - \cost(\bundle{j})$ telescopes exactly
along a path, giving $\pathw{\alloc}(u) = \cost(\bundle{u}) - \min_j
\cost(\bundle{j})$, so it suffices to keep $\max_i \cost(\bundle{i}) - \min_i
\cost(\bundle{i}) \le 1$. That is an invariant of the greedy rule: if all costs
lie in $\set{\mu,\mu+1}$ and the next chore joins a bundle attaining $\mu$, its
cost becomes $\mu$ or $\mu+1$ because marginals lie in $\set{0,1}$, and no other
bundle moves. Observation~\ref{obs:integrality} finishes.
\end{proof}

The case $n = 2$ also holds, with total subsidy at most $1$; there the theorem
is~\cite[Theorem~4]{BKNS22} and ours is only the reduction, which consumes the
identity $\items \setminus \bundle{j} = \bundle{i}$ and so says nothing for
$n \ge 3$.

\begin{theorem}[Small bundles]
\label{main:smallbundle}
Let $\alloc$ be envy-freeable with $\cost_v(\bundle{i}) \le 1$ for all
$v, i \in \agents$. Then $\optsubsidy \in \set{0,1}^n$ and the total subsidy is
at most $n-1$.
\end{theorem}

\begin{proof}
By Theorem~\ref{main:cyclebound}, $\pathw{\alloc}(u) \le \max(0, \max_{v \ne u}
[\cost_v(\bundle{u}) - \cost_v(\bundle{v})]) \le \max(0, 1-0) = 1$.
Observation~\ref{obs:integrality} gives $\optsubsidy \in \set{0,1}^n$, and the
endpoint of a heaviest path has $\pathw{} = 0$.
\end{proof}

\begin{corollary}[Fewer chores than agents]
\label{main:m-le-n}
If $m \le n$ then Conjecture~\ref{conj:main} holds, and a witness is computable
by one minimum-cost bipartite matching.
\end{corollary}

\begin{proof}
Among allocations giving each agent at most one chore take one minimising
$\sum_i \cost_i(\bundle{i})$. Every permutation of its own bundles is an
allocation of the same form, so Theorem~\ref{thm:hs-characterisation}(ii) makes
it envy-freeable; and $\abs{\bundle{i}} \le 1$ forces
$\cost_v(\bundle{i}) \le 1$, so Theorem~\ref{main:smallbundle} applies. The
bound is tight by Example~\ref{ex:lowerbound}, which has $m = n-1$.
\end{proof}

\begin{definition}[Uniformly balanced family]
\label{main:unifbal}
A partition of $\items$ into $n$ bundles is \emph{uniformly balanced} if every
agent values all $n$ bundles within one unit of each other.
\end{definition}

\begin{theorem}[Instances admitting a uniformly balanced family]
\label{main:balanced-class}
Conjecture~\ref{conj:main} holds for every instance admitting a uniformly
balanced family, and a witness is obtained by assigning the bundles by one
maximum-weight matching.
\end{theorem}

\begin{proof}
Assign by a maximum-weight matching; Theorem~\ref{thm:hs-characterisation} makes
the allocation envy-freeable, and uniform balance gives
$\edgew{\alloc}(i,j) \ge -1$ for \emph{any} assignment of that family, so
Corollary~\ref{main:onestep} applies.
\end{proof}

This last case is much the widest of the four. It contains, for instance, every
\emph{symmetric} instance --- one in which each $\cost_i(S)$ depends only on
$\abs{S}$ --- since a partition into bundles of near-equal size shows agent $i$
only the two values $f_i(q)$ and $f_i(q+1)$, which differ by a single marginal.
It is nevertheless not everything: there are instances admitting no uniformly
balanced family at all, and by \Cref{sec:obstructions} no criterion phrased on
arcs alone can reach them.

% --------------------------------------------------------------------------
\subsection{A criterion at three agents}
\label{sec:main-spread}

The fifth result is a criterion rather than a class, and it is the only
statement in this report that bears on three agents in general. For a family
$B = (B_1,\dots,B_n)$ --- a partition of $\items$ into $n$ possibly empty
bundles, not yet assigned to agents --- write
\[
  \mathrm{sp}_i(B) := \max_t \cost_i(B_t) - \min_t \cost_i(B_t),
  \qquad
  \Sigma(B) := \sum_{i \in \agents}\mathrm{sp}_i(B) .
\]
A family is uniformly balanced exactly when every $\mathrm{sp}_i(B) \le 1$.

\begin{theorem}[Small total spread suffices]
\label{main:sigma3}
Let $n = 3$ and let $\sigma$ be any minimum-cost assignment of a family $B$ to
the agents. If an arc of $\envygraph{\alloc}$ has weight at least $2$ then
$\Sigma(B) \ge 4$; and if every arc has weight at most $1$ but some two-path has
weight $2$, then $\Sigma(B) \ge 5$. In particular, if $\Sigma(B) \le 3$ then
\emph{every} minimum-cost assignment of $B$ is good.
\end{theorem}

\begin{proof}
Normalise by $v_i(t) := \cost_i(B_t) - \min_s \cost_i(B_s)$ and put
$x_i := v_i(\sigma(i))$, so that
$\edgew{\alloc}(i,k) = v_i(\sigma(i)) - v_i(\sigma(k))$ and a minimum-cost
assignment is one minimising $F(\sigma) = \sum_i x_i$. If $x_i \ge 1$ then $v_i$
does not vanish at $\sigma(i)$, so its zero lies at some $\sigma(k)$ and the
heaviest arc out of $i$ is exactly $x_i$; if $x_i = 0$ every arc out of $i$ is at
most $0$.

Suppose some arc has weight at least $2$, so some $x_i \ge 2$, and let $k$ hold
the zero of $v_i$. Transposing $i$ and $k$ changes $F$ by
$v_k(\sigma(i)) - x_i - x_k$, non-negative by optimality, so
$v_k(\sigma(i)) \ge x_i + x_k \ge 2$ and $\mathrm{sp}_k \ge 2$; with
$\mathrm{sp}_i \ge x_i \ge 2$ this gives $\Sigma \ge 4$.

Suppose instead every arc is at most $1$ and some two-path $i \to j \to k$ has
weight $2$. Then $x_i = x_j = 1$, $v_i(\sigma(j)) = 0$ and $v_j(\sigma(k)) = 0$,
so $F(\sigma) = 2 + x_k$. Transposing $i$ and $j$ costs
$v_j(\sigma(i)) + x_k$, forcing $v_j(\sigma(i)) \ge 2$ and $\mathrm{sp}_j \ge 2$.
The three-cycle $i \mapsto \sigma(j)$, $j \mapsto \sigma(k)$, $k \mapsto
\sigma(i)$ costs $v_k(\sigma(i))$, forcing $v_k(\sigma(i)) \ge 2 + x_k \ge 2$ and
$\mathrm{sp}_k \ge 2$. With $\mathrm{sp}_i \ge 1$ this gives $\Sigma \ge 5$.

For the last claim, a minimum-cost assignment is envy-freeable by
Theorem~\ref{thm:hs-characterisation}, so by \Cref{prop:n3-main} goodness is
exactly the absence of an arc of weight $2$ and of a two-path of weight $2$, both
of which force $\Sigma \ge 4$.
\end{proof}

\begin{proposition}[Goodness at three agents]
\label{prop:n3-main}
Let $n = 3$ and let $\alloc$ be envy-freeable. Then $\alloc$ is good if and only
if every arc has weight at most $1$ and every directed two-path has weight at
most $1$.
\end{proposition}

\begin{proof}
With no positive cycle, $\pathw{\alloc}(i)$ is the heaviest weight of a
\emph{simple} path leaving $i$; on three vertices such a path has length $0$, $1$
or $2$.
\end{proof}

Theorem~\ref{main:sigma3} strictly contains Theorem~\ref{main:balanced-class} at
$n = 3$: a uniformly balanced family has $\Sigma \le 3$, but so do families with
spread profiles such as $(2,1,0)$ and $(3,0,0)$, which are not uniformly
balanced. It therefore reaches instances admitting no uniformly balanced family
at all --- the ones \Cref{sec:obstructions} shows no arc-based criterion can
certify.

Theorem~\ref{main:sigma3} is conditional on a family with $\Sigma \le 3$
existing. Such a family does always exist for a large class of costs, and that
gives the first case of Conjecture~\ref{conj:main} with three agents.

\begin{definition}[Composed cost]
\label{main:composed}
A cost is \emph{composed} if $\cost_i(S) = f_i(\abs{S \cap D_i})$ for some
$D_i \subseteq \items$ and some monotone $f_i$ with all increments in
$\set{0,1}$.
\end{definition}

Writing $x := \abs{S \cap D_i}$, the composed costs include the binary additive
ones $x$, the capped ones $\min(x,k)$, the threshold ones $\max(0, x-t)$, and
any mixture of these across agents.

\begin{theorem}[Three agents, composed costs]
\label{main:n3composed}
Conjecture~\ref{conj:main} holds for every instance with $n = 3$ in which every
cost is composed.
\end{theorem}

The argument runs as follows; the details are deferred. If the instance admits a
uniformly balanced family, Theorem~\ref{main:balanced-class} already applies. If
it does not, one shows first that some $\abs{D_i}$ must be divisible by $3$, and
then that any three sets admit a $3$-colouring splitting two of them within one
and the third within two. Splitting within one is the same as splitting exactly
evenly for a set whose size is divisible by $3$, so that agent is driven to
spread $0$; a composed cost has spread at most that of its underlying set; and
the three cost spreads come to at most $0 + 1 + 2 = 3$. Theorem~\ref{main:sigma3}
then applies.

\begin{remark}[What remains]
\label{main:rem-spread-gap}
Theorem~\ref{main:n3composed} does not settle $n = 3$ in general. The step from
a set's spread to its cost's spread needs the cost to be a function of the single
quantity $\abs{S \cap D_i}$, and a dichotomous cost need not be of that form; no
reduction of the general case to the composed one is known. For $n \ge 4$
nothing beyond the four cases of \Cref{sec:main-solved} is available, since
Theorem~\ref{main:sigma3} is proved only at three agents --- its proof uses the
fact that a simple path on three vertices has length at most two.
\end{remark}
